\documentclass[11pt]{article}

\usepackage[a4paper,margin=1in]{geometry}
\usepackage{amsmath,amssymb,amsthm,mathtools}
\usepackage{enumitem}
\usepackage[hidelinks]{hyperref}

\newtheorem{theorem}{Theorem}[section]
\newtheorem{proposition}[theorem]{Proposition}
\newtheorem{lemma}[theorem]{Lemma}
\newtheorem{corollary}[theorem]{Corollary}
\theoremstyle{definition}
\newtheorem{remark}[theorem]{Remark}
\newtheorem{problem}[theorem]{Problem}
\newtheorem{definition}[theorem]{Definition}

\newcommand{\R}{\mathbb R}
\newcommand{\Om}{\Omega}
\newcommand{\Fr}{\mathcal A}
\newcommand{\X}{\mathcal X}
\newcommand{\HH}{\mathcal H}
\newcommand{\dT}{\delta_T}
\newcommand{\rstar}{\rho_*}
\newcommand{\rdel}{\rho_\delta}
\newcommand{\wt}{w}
\newcommand{\muT}{\mu}
\newcommand{\eps}{\varepsilon}
\newcommand{\dist}{\operatorname{dist}}
\newcommand{\Per}{P}
\newcommand{\B}{B}
\newcommand{\abs}[1]{\left|#1\right|}
\newcommand{\norm}[1]{\left\|#1\right\|}

\title{Remaining Steps in the Level-Set Rearrangement Route\\
       \large Same-Centre Fusco--Julin, Metric First Variation, and a Bad-Density Repair}
\author{}
\date{May 2026}

\begin{document}
\maketitle

\begin{abstract}
This note records the outcome of attacking the three remaining gaps isolated
in the debug note for the continuous-rearrangement proof.  The first gap,
the same-centre Fusco--Julin input, is actually closed: the correct import is
not the separately minimised statement currently written in
\texttt{preliminaries\_quant\_isoperimetry.tex}, but the centred consequence
of Fusco--Julin's oscillation index together with their Proposition~1.2
(equivalently Proposition~5.5 in Fusco's survey).  The second gap, the
finite-perimeter metric first variation for the rescaled curve
\(\rho^{-1}E_\rho\), is still not proved by the current manuscript: the
Reshetnyak/AGS closure has the lower-semicontinuity inequality in the wrong
direction.  The useful new point is that the third gap, the bad-density
kinetic estimate, can be bypassed if the trace is run on the unscaled
distance
\[
   q(\rho):=\inf_z |E_\rho\Delta B_\rho(z)|
\]
to the moving ball.  Nestedness gives an unconditional Lipschitz bound for
\(q\), so the bad-density set costs \(O(\delta_T)\).  The remaining local
theorem is then a raw moving-ball first-variation estimate on good levels;
in the smooth case it is exactly the vector-field computation below, and in
finite perimeter it is a limsup deformation theorem, strictly weaker than
the old rescaled metric theorem.
\end{abstract}

\tableofcontents

\section{Setup}

Normalize \(|\Omega|=\omega_n\), \(\Omega\subset B_R\), and let \(u\) be the
torsion function.  Write
\[
   E_\rho:=\{u>t(\rho)\},\qquad |E_\rho|=\omega_n\rho^n,
   \qquad \wt(\rho):=-t'(\rho),\qquad d\muT=\wt(\rho)\,d\rho .
\]
The level-set identity gives, on every fixed bulk interval
\([\rstar,\rdel]\subset(0,1]\),
\begin{equation}\label{eq:budget}
   \int_{\rstar}^{\rdel}\bigl(D_H(t(\rho))+D_I(t(\rho))\bigr)\,d\muT(\rho)
   \le C(n,\rstar)\,\dT(\Omega).
\end{equation}
Here \(D_H\) is the Cauchy--Schwarz defect of \(|\nabla u|\) on
\(\partial^*E_\rho\), and
\[
   D_I(t(\rho))=P(E_\rho)^2-P(B_\rho)^2 .
\]
Talenti's profile-gap estimate gives the bad-density bound
\begin{equation}\label{eq:bad-density}
   |B_\tau|\le C(n,\rstar)\,\dT(\Omega),\qquad
   B_\tau:=\{\rho\in[\rstar,\rdel]:\wt(\rho)<\tau_0\},\quad
   \tau_0:=\frac{\rstar}{2n}.
\end{equation}

\section{Same-Centre Fusco--Julin Is Closed}

\subsection{Where the current repository proof goes wrong}

The proof of the scaled same-centre theorem in
\texttt{Manuscript/preliminaries\_quant\_isoperimetry.tex} currently chooses
an oscillation-minimising centre \(z\) and then uses
\[
   \frac{|E\Delta B_\rho(z)|}{\omega_n\rho^n}\ge \Fr(E)
\]
as if it were the needed upper bound.  This is the precise sign error: the
Fraenkel asymmetry is an infimum over centres, so it gives only a lower bound
on the centred symmetric difference at a prescribed centre.  That proof
cannot be patched by algebra.

The statement, however, is available from the correct Fusco--Julin input.
Fusco--Julin do not merely give two unrelated minimisations.  Their
oscillation index is minimised over balls, and their Poincare-type estimate
shows that the same ball controls the \(L^1\) displacement.

\subsection{Import-ready statement}

For \(|E|=|B_\rho|\), define the pointed oscillation
\[
   \beta(E,z)^2
   :=
   \frac12\int_{\partial^*E}
   \left|\nu_E(x)-\frac{x-z}{|x-z|}\right|^2\,d\HH^{n-1}(x).
\]

\begin{theorem}[Same-centre Fusco--Julin package]\label{thm:scFJ}
Let \(n\ge2\) and \(0<\rstar\le\rho\le1\).  There is a constant
\(C_{\rm scFJ}=C_{\rm scFJ}(n,\rstar)\), explicitly traceable to the
Fusco--Julin constants, with the following property.  If
\(E\) is a finite-perimeter set with \(|E|=\omega_n\rho^n\), then there is a
centre \(z_E\) such that
\begin{equation}\label{eq:scFJ}
   |E\Delta B_\rho(z_E)|^2
   +
   \int_{\partial^*E}
   \left|\nu_E(x)-\frac{x-z_E}{|x-z_E|}\right|^2\,d\HH^{n-1}(x)
   \le
   C_{\rm scFJ}\,\bigl(P(E)-P(B_\rho)\bigr).
\end{equation}
Consequently, on a level \(E_\rho\),
\begin{equation}\label{eq:scFJ-DI}
   |E_\rho\Delta B_\rho(z_\rho)|^2
   +
   \int_{\partial^*E_\rho}
   \left|\nu_\rho-\frac{x-z_\rho}{|x-z_\rho|}\right|^2\,d\HH^{n-1}
   \le C(n,\rstar)D_I(t(\rho)).
\end{equation}
\end{theorem}

\begin{proof}
After translating and scaling, it is enough to treat \(\rho=1\).  Let \(z_E\)
be a minimiser of the Fusco--Julin oscillation index.  Fusco--Julin's strong
form gives
\[
   \beta(E,z_E)^2\le C_{FJ}(n)\bigl(P(E)-P(B_1)\bigr).
\]
Their Poincare-type estimate for the same minimising ball gives, at this
same centre,
\[
   |E\Delta B_1(z_E)|^2
   \le C(n)\,\beta(E,z_E)^2 .
\]
Equivalently, in the notation of \cite[Prop.~1.2]{FJ2014} or
\cite[Prop.~5.5]{FuscoSurvey2015}, once the oscillation centre has been
translated to the origin,
\[
   \beta(E,0)^2\ge c(n)|E\Delta B_1|^2
\]
up to the harmless perimeter-deficit term that is already dominated by the
same strong inequality.  Combining the two estimates proves
\eqref{eq:scFJ} for \(\rho=1\).

For general \(\rho\), apply the unit-radius statement to
\(\widetilde E:=\rho^{-1}(E-z_E)\).  The oscillation integral scales like
\(\rho^{1-n}\), the perimeter deficit like \(\rho^{1-n}\), and
\(|E\Delta B_\rho(z_E)|^2=\rho^{2n}|\widetilde E\Delta B_1|^2\).  Since
\(\rho\ge\rstar\), all scale factors are absorbed into
\(C_{\rm scFJ}(n,\rstar)\).  Finally,
\[
   P(E_\rho)-P(B_\rho)
   =
   \frac{D_I(t(\rho))}{P(E_\rho)+P(B_\rho)}
   \le C(n,\rstar)D_I(t(\rho)),
\]
which gives \eqref{eq:scFJ-DI}.
\end{proof}

\begin{remark}[Effectivity of the constant]
The argument above introduces no compactness beyond the constant already
present in the Fusco--Julin theorem.  Thus the same-centre repair is
quantitative in exactly the same sense as the imported Fusco--Julin
inequality: once the constants in Theorem~5.4 and Proposition~5.5 of the
source are unpacked, \(C_{\rm scFJ}\) is obtained by algebra and scaling.
\end{remark}

\begin{remark}
This also repairs the measurable centre selection: choose the minimisers of
the joint closed sublevel set supplied by \eqref{eq:scFJ}.  The selection is
standard once the centre is restricted to a compact ball; compactness follows
from \(E_\rho\subset B_R\) and the centred \(L^1\) bound.
\end{remark}

\section{The Current Rescaled Metric Closure Is Not Sound}

The manuscript file \texttt{Manuscript/plan2\_metric\_framework.tex} tries to
pass from smooth approximations to finite-perimeter level sets by combining
Reshetnyak lower semicontinuity with AGS lower semicontinuity of metric
derivatives.  The decisive lines have the following form.  For smooth
approximants,
\[
   |\dot F^k_\rho|_\X\le m_k(\rho),
\]
and weak compactness gives \(m_k\rightharpoonup m\).  Reshetnyak gives only
\[
   m(\rho)\ge
   \rho_*^{-n}\inf_a
   \int_{\partial^*E_\rho}
   |V_\rho-H_{z_\rho,\rho}-a\cdot\nu_\rho|\,d\HH^{n-1}.
\]
AGS then gives \(|\dot F_\rho|_\X\le m(\rho)\).  These two inequalities do
not imply the desired estimate by the boundary action; one would need
\(m(\rho)\le\) that action, or a direct limsup recovery.  Thus the old
rescaled finite-perimeter metric theorem remains open.

\section{A Bad-Density Bypass: Use the Moving Ball}

The bad-density obstruction was caused by the rescaled curve
\[
   F_\rho=\rho^{-1}E_\rho .
\]
The rescaling differentiates rough sets, and boundedness plus nestedness do
not imply a perimeter-free Lipschitz bound for
\(\rho\mapsto \rho^{-1}E_\rho\).  The following scalar quantity avoids that
operation:
\begin{equation}\label{eq:q-def}
   q(\rho):=\inf_{z\in\R^n}|E_\rho\Delta B_\rho(z)|.
\end{equation}

\begin{lemma}[Unconditional Lipschitz bound for \(q\)]\label{lem:qLip}
For \(0<\rho\le\sigma\le1\),
\[
   |q(\sigma)-q(\rho)|
   \le 2\omega_n(\sigma^n-\rho^n)
   \le 2n\omega_n|\sigma-\rho|.
\]
In particular \(q\in W^{1,\infty}_{\rm loc}((0,1))\), with a computable
constant depending only on \(n\).
\end{lemma}

\begin{proof}
Fix \(z\) with
\(|E_\rho\Delta B_\rho(z)|\le q(\rho)+\eps\).  Since the level sets are
nested and the same-centre balls are nested,
\[
\begin{aligned}
   |E_\sigma\Delta B_\sigma(z)|
   &\le |E_\rho\Delta B_\rho(z)|
      +|E_\sigma\setminus E_\rho|+|B_\sigma(z)\setminus B_\rho(z)|  \\
   &= |E_\rho\Delta B_\rho(z)|
      +2\omega_n(\sigma^n-\rho^n).
\end{aligned}
\]
Taking the infimum in \(z\) and then \(\eps\downarrow0\) gives
\[
   q(\sigma)\le q(\rho)+2\omega_n(\sigma^n-\rho^n).
\]
Interchanging \(\rho\) and \(\sigma\) gives the reverse inequality.
\end{proof}

\begin{corollary}[Bad-density action is free for \(q\)]\label{cor:badtau}
With \(B_\tau\) as in \eqref{eq:bad-density},
\[
   \int_{B_\tau}|q'(\rho)|^2\,d\rho
   \le C(n,\rstar)\,\dT(\Omega).
\]
\end{corollary}

\begin{proof}
By Lemma~\ref{lem:qLip}, \(|q'|\le 2n\omega_n\) a.e.  Combine this with
\eqref{eq:bad-density}.
\end{proof}

\section{The Scalar Trace Lemma}

The moving-ball distance gives the endpoint extraction without any rescaled
bad-level kinetic estimate.

\begin{lemma}[Scalar moving-ball trace]\label{lem:scalartrace}
Let \(J=[a,b]\subset[\rstar,\rdel]\) have length bounded below by
\(\ell_0>0\).  Assume
\begin{align}
   &\int_J q(\rho)^2\,d\muT(\rho)\le C_0\dT, \label{eq:traceH1}\\
   &\int_{J\cap G_\tau}|q'(\rho)|^2\,d\rho\le C_1\dT, \qquad
     G_\tau:=J\setminus B_\tau, \label{eq:traceH2}\\
   &|B_\tau\cap J|\le C_2\dT . \label{eq:traceH3}
\end{align}
Then, for \(\dT\) small enough depending on \(n,\rstar,\ell_0\),
\[
   \sup_{\rho\in J}q(\rho)^2\le C(n,\rstar,\ell_0,C_0,C_1,C_2)\,\dT .
\]
\end{lemma}

\begin{proof}
First convert the weighted \(L^2\) distance into a Lebesgue \(L^2\) distance.
On \(G_\tau\), \(d\muT\ge\tau_0\,d\rho\), hence
\[
   \int_{J\cap G_\tau}q^2\,d\rho\le \tau_0^{-1}C_0\dT .
\]
On \(B_\tau\), the trivial bound \(q(\rho)\le |E_\rho|+|B_\rho|=2\omega_n\rho^n
\le2\omega_n\) and \eqref{eq:traceH3} give
\[
   \int_{J\cap B_\tau}q^2\,d\rho\le 4\omega_n^2 C_2\dT.
\]
Thus \(\int_J q^2\,d\rho\le C\dT\), and there is \(\rho_0\in J\) with
\[
   q(\rho_0)^2\le \frac{C}{|J|}\dT\le C(\ell_0)\dT.
\]
Next,
\[
   \int_J |q'|^2\,d\rho
   =
   \int_{J\cap G_\tau}|q'|^2\,d\rho+
   \int_{J\cap B_\tau}|q'|^2\,d\rho
   \le C\dT
\]
by \eqref{eq:traceH2} and Corollary~\ref{cor:badtau}.  For every
\(\rho\in J\),
\[
   q(\rho)
   \le q(\rho_0)+\int_J|q'|
   \le C\dT^{1/2}+|J|^{1/2}\left(\int_J|q'|^2\right)^{1/2}
   \le C\dT^{1/2}.
\]
Squaring proves the claim.
\end{proof}

\section{The Remaining Local Theorem}

The preceding sections reduce the metric part to a local good-level estimate
for \(q\).  This theorem is the precise replacement for the old rescaled
finite-perimeter metric theorem.

\begin{problem}[Raw moving-ball first variation]\label{prob:rawFV}
Let \(E_\rho=\{u>t(\rho)\}\) be a.e. finite-perimeter levels.  On the good
set where \(D_I(t(\rho))\) is small, choose the same-centre
Fusco--Julin centre \(z_\rho\) from Theorem~\ref{thm:scFJ}.  Prove the
finite-perimeter limsup estimate
\begin{equation}\label{eq:rawFV}
   |q'(\rho)|
   \le C(n,R,\rstar)q(\rho)
   +C(n,R,\rstar)\inf_{a\in\R^n}
      \int_{\partial^*E_\rho}
      |V_\rho-e_{z_\rho}\cdot\nu_\rho-a\cdot\nu_\rho|
      \,d\HH^{n-1}
\end{equation}
for a.e. good \(\rho\), where
\[
   V_\rho=\frac{\wt(\rho)}{|\nabla u|},\qquad
   e_{z_\rho}(x)=\frac{x-z_\rho}{|x-z_\rho|}.
\]
\end{problem}

\begin{proposition}[Why \eqref{eq:rawFV} closes the route]\label{prop:rawcloses}
If Problem~\ref{prob:rawFV} is proved, then the bounded sharp
Saint--Venant estimate
\[
   \Fr(\Omega)^2\le C(n,R)\dT(\Omega)
\]
follows from the level-set identity and the boundary-layer transfer, with
computable constants.
\end{proposition}

\begin{proof}
On a good level, Theorem~\ref{thm:scFJ} gives
\[
   q(\rho)^2\le C D_I(t(\rho)),
   \qquad
   \int_{\partial^*E_\rho}|1-e_{z_\rho}\cdot\nu_\rho|\,d\HH^{n-1}
   \le C D_I(t(\rho)).
\]
The velocity variance identity gives
\[
   \left(\int_{\partial^*E_\rho}|V_\rho-\bar V_\rho|\,d\HH^{n-1}\right)^2
   \le C D_H(t(\rho)),
   \qquad
   \bar V_\rho:=\frac{P(B_\rho)}{P(E_\rho)} .
\]
Since \(P(E_\rho)-P(B_\rho)\le C D_I(t(\rho))\) on the good set,
\[
\begin{aligned}
   \int_{\partial^*E_\rho}|V_\rho-e_{z_\rho}\cdot\nu_\rho|
   &\le
   \int|V_\rho-\bar V_\rho|
   +P(E_\rho)|\bar V_\rho-1|
   +\int|1-e_{z_\rho}\cdot\nu_\rho|                                      \\
   &\le C D_H(t(\rho))^{1/2}+C D_I(t(\rho)).
\end{aligned}
\]
Squaring and using \(D_I\le\eta_I\) on the good set,
\[
   |q'(\rho)|^2\le C\bigl(D_H(t(\rho))+D_I(t(\rho))\bigr)
\]
there.  On \(G_\tau\), \(d\rho\le\tau_0^{-1}d\muT\), so
\[
   \int_{G_\tau}|q'|^2\,d\rho
   \le C\int_{\rstar}^{\rdel}(D_H+D_I)\,d\muT
   \le C\dT
\]
by \eqref{eq:budget}.  The distance hypothesis
\(\int q^2\,d\muT\le C\dT\) follows from \(q^2\le C D_I\) on the good
isoperimetric set and the usual bad-\(I\) weighted Markov estimate with the
trivial bound \(q\le2\omega_n\).  Lemma~\ref{lem:scalartrace} gives
\[
   q(\rdel)^2\le C\dT
\]
on an order-one window ending at \(\rdel\) (or at the good endpoint
\(\widehat\rho\)).  Therefore
\[
   \Fr(E_{\rdel})^2
   \le C(n,\rstar)q(\rdel)^2
   \le C\dT .
\]
Finally \(E_{\rdel}\subset\Omega\) and
\[
   |\Omega\setminus E_{\rdel}|=\omega_n(1-\rdel^n)
   \le C(n)\sqrt{\dT},
\]
so the elementary boundary-layer transfer gives
\[
   \Fr(\Omega)\le \Fr(E_{\rdel})+C\sqrt{\dT},
\]
and the desired squared estimate follows.
\end{proof}

\subsection{Smooth proof of the local theorem}

For smooth moving level surfaces, Problem~\ref{prob:rawFV} is exactly the
standard relative first variation.  Fix a centre \(z\) and a translation
velocity \(a\).  Let \(X=e_z+a\) near \(\partial B_\rho(z)\), with a harmless
smooth cutoff away from that sphere so that the flow is globally Lipschitz.
The flow of \(B_\rho(z)\) has boundary normal velocity \(1+a\cdot\nu\), hence
it agrees to first order with \(B_{\rho+h}(z+ha)\).  The level surface
\(\partial E_\rho\) has normal velocity \(V_\rho\).  Differentiating
\(|E_\rho\Delta B_\rho(z)|\) under these two motions gives
\[
   \frac{d^+}{d\rho}|E_\rho\Delta B_\rho(z)|
   \le
   C q(\rho)
   +
   \int_{\partial E_\rho}
      |V_\rho-e_z\cdot\nu_\rho-a\cdot\nu_\rho|\,d\HH^{n-1}.
\]
The \(Cq\) term is the Jacobian error on the already mismatched region and is
absorbed by \(q^2\le C D_I\).  Taking the infimum over \(z\) and \(a\), and
then using measurable near-minimisers, gives \eqref{eq:rawFV}.

The finite-perimeter version should be proved by the same spacetime
Gauss--Green argument applied to
\[
   \mathcal E:=\{(\rho,x):x\in E_\rho\}\subset[\rstar,\rdel]\times B_R,
\]
whose perimeter is supplied by the coarea formula.  This is a genuine limsup
deformation theorem, not a lower-semicontinuity passage.  It is the remaining
load-bearing theorem.

\section{Verdict}

The state after this attack is sharper than before.

\begin{enumerate}[leftmargin=2em]
\item The same-centre Fusco--Julin input is closed, but the proof in the
current preliminaries should be replaced.  The correct source is
Fusco--Julin's same oscillation centre plus their Poincare-type control of
the symmetric difference.
\item The rescaled finite-perimeter metric theorem is still not proved by
the manuscript's Reshetnyak/AGS passage; the inequality direction is wrong.
\item The bad-density kinetic obstruction is not intrinsic.  It is an
artifact of rescaling rough sets.  For the unscaled moving-ball distance
\(q(\rho)\), nestedness gives a uniform Lipschitz bound and the bad-density
set costs \(O(\dT)\).
\item The remaining theorem is now the raw moving-ball first variation
\eqref{eq:rawFV}.  It is local, good-level, and finite-perimeter; proving it
with a direct spacetime limsup deformation would close the bounded
Saint--Venant route without invoking BDV stability.
\end{enumerate}

\begin{thebibliography}{9}

\bibitem{FJ2014}
N.~Fusco and V.~Julin,
\emph{A strong form of the quantitative isoperimetric inequality},
Calc. Var. Partial Differential Equations \textbf{50} (2014), 925--937,
\url{https://arxiv.org/abs/1111.4866}.

\bibitem{FuscoSurvey2015}
N.~Fusco,
\emph{The quantitative isoperimetric inequality and related topics},
Bull. Math. Sci. \textbf{5} (2015), 517--607,
\url{https://link.springer.com/article/10.1007/s13373-015-0074-x}.

\bibitem{Maggi2012}
F.~Maggi,
\emph{Sets of Finite Perimeter and Geometric Variational Problems},
Cambridge University Press, 2012.

\bibitem{AFP2000}
L.~Ambrosio, N.~Fusco, and D.~Pallara,
\emph{Functions of Bounded Variation and Free Discontinuity Problems},
Oxford University Press, 2000.

\end{thebibliography}

\end{document}
